\section{What the certificate does and does not transport}\label{sec:discussion}

The theorem joins five representations of the same finite object:
\[
\begin{aligned}
 \text{link Haar variables}
 &\longrightarrow \text{tree-gauge quotient}
 \longrightarrow \text{logarithmic curvature chart}\\
 &\longrightarrow \text{retained/fiber coordinates}
 \longrightarrow \text{effective retained action}.
\end{aligned}
\]
At each arrow, a different invariant is preserved.  Tree gauge preserves the
quotient integral exactly.  The logarithmic chart preserves the exact
plaquette action while making nonlinearity quantitative.  The
retained--fiber splitting preserves the tangent quadratic form without a
mixed term.  Marginalization preserves normalization and yields a controlled
effective Hessian, although it forgets the eliminated fiber coordinate.

This is the useful content of the phrase ``one-block certificate.''  It is
not a small mass gap.  The floor $1/25$ is a coordinate-dependent eigenvalue
of a fixed finite matrix.  It carries neither physical units that survive a
lattice-spacing limit nor a statement about the spectrum of a reconstructed
Hamiltonian.

The theorem also identifies its own failure mode.  The local nonlinear
inequalities apply only in $\cC_r$.  Proposition~\ref{prop:reference-tail}
shows that the cylinder is not vacuous for the tangent Gaussian at the chosen
time, but the interacting density can redistribute normalization through
shared plaquettes, boundary conditions, large fields, cut-locus regions, and
topological sectors.  A valid next theorem must control those effects in a
dependent, gauge-invariant complement norm.  Replacing $10^{-18}$ by a more
attractive decimal without such a theorem would not improve the logical
position.

If a complement estimate were available, two additional layers would still
be required.  First, the generated action must remain in a quasilocal class
under repeated blocking, with constants uniform in the scale and volume.
Second, a continuum construction must preserve reflection positivity and a
dense gauge-invariant observable algebra strongly enough to identify the
Hamiltonian spectrum.  None of these statements is an implicit consequence
of \eqref{eq:effective-coercivity}.

The certificate is therefore best viewed as a reproducible interface lemma.
It says precisely what a future large-field or multiscale argument may import,
and precisely which information that argument must supply independently.
